\section{What was verified by machine}\label{sec:verification}

Apart from the statements transcribed in \cref{sec:statements}, everything in this section was
computed by a program run on the archive. The programs and their output are public
(\cref{app:reproduce}).

\subsection{The archive}
The Lean source is distributed as fourteen parts which concatenate to a single file
\lean{lean.tar.gz} of \archivebytes{} bytes with SHA-256
\[
\text{\small\ttfamily 9631b4f020199e580ba86178cb56c6e9}\;\text{\small\ttfamily b09f690c2369a80dc4334aa93f49c7fc}.
\]
It unpacks to \filecount{} \lean{.lean} files (\linecount{} lines) and three build files. The
build files pin Lean \leanversion{} and Mathlib at commit
\lean{0df444a360eaa60ab8c11dca51a86af692955474}.

\subsection{Static inspection}
Before building we scanned the text of every file, comments included. No file contains the
word \lean{sorry} or \lean{admit}, an \lean{axiom} declaration, or any of
\lean{native_decide}, \lean{implemented_by}, \lean{extern}, \lean{unsafe}, \lean{opaque},
\lean{csimp}, or the names by which a declaration can be added to an environment without
passing the kernel. Every import resolves to a file of the archive, of Mathlib or of Lean; the
import graph is acyclic; every file is in the import closure of \lean{Main.lean}. Beside 87
macros, which abbreviate syntax, the text defines four elaborators in three files: two
tactics, a command that copies declarations under new names, and a probe. Each produces its
result through an ordinary declaration, so the kernel checks what they produce. Five small
files run a command while they are compiled; it adds a name to a deny list of Lean's
suggestion facility and declares nothing. Two helper functions of the elaborators are declared
\lean{partial}; they are programs, not statements, and they are outside the dependency cone of
\cref{sec:cone}.

\subsection{The statements}\label{sec:statements}
\lean{Main.lean} is twelve lines long and proves five theorems.
\begin{statement}[\lean{riemannHypothesis}]\label{st:rh}
For every $s\in\C$ with $\zeta(s)=0$, if $s$ is not of the form $-2(n+1)$ with $n\in\N$ and
$s\neq 1$, then $\re s=\tfrac12$.
\end{statement}
The statement is the proposition \lean{RiemannHypothesis} of Mathlib, not a definition of the
archive: in the compiled environment that name is declared in the module
\lean{Mathlib.NumberTheory.LSeries.RiemannZeta}, as is \lean{riemannZeta}.
\begin{statement}[\lean{LFunction_ne_zero_of_half_lt_re}]\label{st:grh}
For every modulus $N\ge 1$, every Dirichlet character $\chi$ modulo $N$ and every $s\in\C$ with
$\re s>\tfrac12$, if $\chi\neq 1$ or $s\neq1$ then $L(s,\chi)\neq0$.
\end{statement}
\begin{statement}[\lean{LFunction_eq_zero_re}]\label{st:strip}
For every modulus $N\ge1$, every Dirichlet character $\chi$ modulo $N$ and every $s$ with
$0<\re s<1$ and $L(s,\chi)=0$, one has $\re s=\tfrac12$.
\end{statement}
The remaining two are \lean{riemannZeta_eq_zero_re}, which is \cref{st:rh} again, and
\lean{LFunction_eq_zero_re_of_isPrimitive}, which is \cref{st:strip} for primitive characters.
Here $L(s,\chi)$ is Mathlib's \lean{DirichletCharacter.LFunction}, the meromorphic continuation.
The characters are not assumed primitive; restricting to the open strip in \cref{st:strip}
excludes the zeros on $\re s=0$ that imprimitive characters have, so the statements are
correctly posed.

``GRH'' in this paper always means \cref{st:grh,st:strip}: Dirichlet $L$-functions. The note
says the case of Dedekind zeta functions is expected and not done.

\subsection{The build}
We ran \lean{lake build Main} on the unmodified archive, on a rented machine with 32 cores,
using the Mathlib build cache for the pinned commit. It finished after \buildhours{} hours
with the line \lean{Build completed successfully (183893 jobs)} and exit code~$0$. All
\filecount{} modules of the archive were built; none failed.

\subsection{Axioms}
For each of the five theorems, \lean{#print axioms} reports exactly
\[
\text{\lean{propext}},\qquad \text{\lean{Classical.choice}},\qquad \text{\lean{Quot.sound}},
\]
the three axioms on which Mathlib as a whole rests. In particular none depends on
\lean{sorryAx} or on an axiom that trusts the compiler.

\subsection{The dependency cone}\label{sec:cone}
A build proves that every file compiles; it does not say which declarations the final theorem
uses. From the compiled environment we computed the set of all constants reachable from
\lean{riemannHypothesis} through the type and the value of every declaration reached. It has
\conesize{} elements. Of the \projtheorems{} theorems the archive declares,
\coneprojtheorems{} are in it, and the cone meets \coneprojmodules{} of the \projmodules{}
modules that declare anything. These are counts of names: Lean lets an auxiliary theorem be
declared again in each module that needs it, and \repeatedtheorems{} of the archive's theorem
declarations repeat a name declared in another module. When this paper says that a declaration
is, or is not, used by the proof, it means membership in this set. \Cref{sec:offpath} lists
what is outside it.

\subsection{Kernel replay}
A build checks each declaration once, in the process that elaborated it. A replay passes the
declarations to the kernel a second time, from the compiled files, in a process that does
nothing else. We replayed the whole import closure of \lean{Main.lean} into an environment that
starts empty: the archive, Mathlib with its dependencies, and Lean's own library,
\replaymodules{} modules with \replayconstants{} declarations. Of these, \replaydecls{} are
neither unsafe nor partial, and they are what a replay checks. Each of them, other than a
constructor or a recursor, is passed to the kernel's entry point
\lean{Kernel.Environment.addDeclCore} against an environment containing only declarations that
precede it in a dependency order; constructors and recursors are compared with the ones the
kernel generates from the inductive types.

\replaystatus

The toolchain's checker, \lean{leanchecker} \cite{lean4checker}, cannot do this on an archive of
this shape. It imports the whole import closure afresh for every module it checks, and its
first batch of 2{,}012 modules exceeded the machine's memory limit of 128\,GB; its other mode
replays in one process but on one thread. The replay was therefore done by a program of about
three hundred lines, published with this paper. It imports once, orders the declarations with
a traversal that is linear in the size of each term, and runs the kernel checks in parallel,
each against a snapshot of the environment as it stood when the declaration's turn came;
inductive types are added on the main thread, as in Lean's own \lean{Environment.replay}. The
run ends with a test of the checker itself: it must accept a proof of \lean{True} and reject
the same term offered as a proof of \lean{False}.

\subsection{What this establishes, and what it does not}
In Lean \leanversion{} with Mathlib at the pinned commit, \cref{st:rh,st:grh,st:strip} are
theorems, with no assumption beyond the standard axioms. That is the whole of what the machine
certifies. It does not certify any description of the proof, including ours.

Four limits should be stated. Mathlib was not re-elaborated from source: its compiled files
came from the community cache for the pinned commit, and what we add is that the kernel
re-checked its declarations in the replay. The replay program is ours and has not been reviewed
by anyone else; its self-test shows only that it rejects one false proof, and the build's own
kernel checks do not depend on it. We know of no other build of the archive at the time of
writing, and an independent one would be worth having. And the soundness of Lean's kernel is
assumed, as in every formalisation.
